\chapter{Computational Modelling of Reactor Dynamics}
\label{ch:comp}

The theory of the preceding chapters becomes vivid when integrated numerically.
This chapter describes a small, open-source point-reactor-kinetics toolkit written
to accompany this book, shows the transients it produces, and uses it to make the
stability principles concrete---including a numerical caricature of the Chernobyl
feedback failure.

\section{A minimal point-kinetics solver}

The companion code, \texttt{reactor-kinetics}, couples validated delayed-neutron
data to a stiff ODE integrator with optional temperature feedback. Its core is the
six-group point kinetics of Chapter~\ref{ch:pk} together with the adiabatic
feedback model of Chapters~\ref{ch:thermal} and~\ref{ch:nonlinear}:
\begin{align*}
\frac{\dd p}{\dd t} &= \frac{\rho(t)-\beta}{\Lambda}p + \sum_i\lambda_i\zeta_i,\quad
\frac{\dd \zeta_i}{\dd t} = \frac{\beta_i}{\Lambda}p - \lambda_i\zeta_i,\\
\frac{\dd T}{\dd t} &= \frac{P_0\,p}{C},\qquad
\rho(t) = \rho_{ext}(t) + \alpha\,(T-T_0).
\end{align*}
Because $\Lambda\sim 10^{-5}\unit{s}$ while the precursor times are seconds, the
system is stiff; the solver uses an implicit/adaptive method (LSODA). Delayed-neutron
parameters are taken from validated $^{235}$U data (Appendix~\ref{app:data}).

\section{Verification against the inhour equation}

A dynamics code is only trustworthy if it reproduces known analytical results. The
toolkit's asymptotic period, measured from simulated power traces, matches the
inhour equation \eqref{eq:inhour} to better than $10^{-4}$ across a range of step
reactivities (Figure~\ref{fig:inhour}). This agreement validates both the solver
and the kinetics implementation, and it reproduces the steep collapse of the period
as reactivity approaches one dollar---the quantitative warning sign of approaching
prompt criticality.

\section{A gallery of transients}

Figure~\ref{fig:transients} collects four canonical transients computed with the
toolkit, each illustrating a principle from earlier chapters:

\begin{enumerate}[nosep]
\item \textbf{Sub-prompt step ($+0.3\$$).} A prompt jump
(Eq.~\ref{eq:promptjump}) followed by a slow rise on a stable delayed-neutron
period---the controllable regime.
\item \textbf{Prompt-critical step ($+1.2\$$).} Runaway growth on the prompt
timescale (note the logarithmic axis), the signature of crossing one dollar.
\item \textbf{Control-rod scram.} A negative ramp drives the power down, with the
characteristic delayed-neutron tail as the precursors decay.
\item \textbf{Reactivity-initiated excursion with feedback.} A prompt-critical
insertion limited by negative temperature feedback: the power peaks and turns
over, the Nordheim--Fuchs self-limitation of Chapter~\ref{ch:nonlinear}.
\end{enumerate}

The fourth panel is the crux. With negative feedback the excursion is capped and
the reactor survives. The exercise below shows what happens when the feedback sign
is wrong.

\section{A numerical caricature of the Chernobyl feedback failure}

The toolkit makes it easy to flip the sign of the feedback and watch stability
turn into instability. Consider two runs with identical prompt-critical insertions
($\rho_0\approx 1.05\$$):

\begin{itemize}[nosep]
\item \textbf{Negative feedback} ($\alpha<0$, the well-designed reactor): the
fuel heats, reactivity falls, and the power excursion turns over at a finite peak,
depositing a bounded energy $E_{\text{total}}\approx 2(\rho_0-\beta)/\gamma$
(Chapter~\ref{ch:nonlinear}). The reactor self-limits.
\item \textbf{Positive feedback} ($\alpha>0$, the Chernobyl-like configuration):
the rising power adds further reactivity, there is no turnover, and the power grows
until the model's adiabatic assumptions break---representing, in reality, fuel
fragmentation, coolant flashing, and disassembly.
\end{itemize}

This two-line change in a parameter reproduces, in caricature, the essential
physics of the accident: the difference between Chernobyl and a survivable RIA is
the \emph{sign of} $\alpha$. The toolkit also includes a model void/temperature
feedback that becomes positive at low ``power,'' qualitatively mimicking the RBMK's
low-power instability window of Chapter~\ref{ch:linear}.

\begin{remark}[On responsible modelling]
The companion code is an educational point-kinetics tool. It is deliberately
\emph{not} a high-fidelity simulation of any specific reactor: it uses lumped
parameters, an adiabatic feedback, and generic data. Its value is pedagogical---to
let a reader reproduce the inhour equation, watch a prompt-critical excursion, and
see negative feedback cap it---not to model a real plant, which requires validated,
qualified codes (e.g.\ coupled neutronics/thermal-hydraulics systems) far beyond a
point model.
\end{remark}

\section{Toward higher fidelity}

Real safety analysis goes well beyond point kinetics. Spatial effects
(Chapter~\ref{ch:xenon}) demand at least nodal or full three-dimensional
neutronics coupled to a multi-channel thermal-hydraulics solver; the RBMK's size
made such spatial modelling essential, and core-average point parameters
understated the danger. Modern coupled codes solve the time-dependent
neutron-transport (or diffusion) equation on a spatial mesh, feed local powers to a
thermal-hydraulic subchannel model, and return local temperatures and densities as
feedback---closing the loop of Chapter~\ref{ch:coupled} cell by cell. The point
model of this chapter is the conceptual seed of those codes, and the place to build
intuition before trusting their output.

\keypoint{A small point-kinetics code reproduces the inhour equation to $10^{-4}$,
displays the canonical transients, and---by flipping the sign of the feedback
coefficient---turns a self-limiting excursion into a Chernobyl-like runaway. The
sign of the feedback is, numerically as well as physically, the whole story.}



\section{Stiffness and absolute stability}
\label{sec:stiffness}

The point-kinetics generator $\mathbf M$ has eigenvalues spanning
$|s_0|\sim\lambda_i\sim10^{-2}$--$10^{0}\unit{s^{-1}}$ up to
$|s_G|\sim\beta/\Lambda\sim10^{3}\unit{s^{-1}}$, a \emph{stiffness ratio}
\begin{equation}
S=\frac{\max_j|\Re s_j|}{\min_j|\Re s_j|}\sim\frac{\beta/\Lambda}{\lambda_1}
\sim10^{5}\text{--}10^{6}.
\end{equation}
An explicit method applied to a mode $\dot y=s y$ is stable only if $hs$ lies in
its bounded region of absolute stability; for forward Euler this requires
$h<2/|s_G|\sim2\Lambda\sim10^{-5}\unit{s}$, prohibitively small. Implicit methods
remove this restriction.

\begin{definition}[A- and L-stability]
A one-step method with stability function $R(z)$ (the amplification for
$\dot y=sy$, $z=hs$) is \emph{A-stable} if its stability region contains the entire
left half-plane, $|R(z)|\le1$ for all $\Re z\le0$. It is \emph{L-stable} if in
addition $R(z)\to0$ as $\Re z\to-\infty$ (rapid damping of the stiffest modes).
\end{definition}

\begin{theorem}[$\theta$-method stability]
\label{thm:theta}
The $\theta$-method $y_{n+1}=y_n+h[(1-\theta)f_n+\theta f_{n+1}]$ has
$R(z)=\dfrac{1+(1-\theta)z}{1-\theta z}$. It is A-stable iff $\theta\ge\tfrac12$;
the trapezoidal rule $\theta=\tfrac12$ is A-stable but not L-stable
($R(-\infty)=-1$), while backward Euler $\theta=1$ is L-stable
($R(-\infty)=0$).
\end{theorem}
\begin{proof}
Apply the method to $\dot y=sy$: $y_{n+1}=R(z)y_n$ with the stated $R$. For
$\Re z\le0$, $|R(z)|\le1\iff|1+(1-\theta)z|\le|1-\theta z|$; squaring and
simplifying gives $(1-2\theta)|z|^2\le -2\Re z\,(\,\cdot\,)$ which holds for all
$\Re z\le0$ iff $1-2\theta\le0$, i.e.\ $\theta\ge\tfrac12$. The limits
$R(-\infty)=(\theta-1)/\theta$ give $-1$ at $\theta=\tfrac12$ and $0$ at
$\theta=1$.
\end{proof}

The backward-differentiation formulas (BDF$k$) used by production solvers (e.g.\
LSODA, CVODE) are $A(\alpha)$-stable up to order $6$ and L-stable at orders $1$--$2$;
their use is what makes the figures of this book reproducible at large step sizes.

\section{Exponential and operator-splitting integrators}
\label{sec:expint}

Because the point-kinetics block is \emph{linear} in $\mathbf y$ for fixed $\rho$,
the exact one-step propagator over an interval of constant reactivity is the matrix
exponential
\begin{equation}
\mathbf y_{n+1}=e^{\mathbf M(\rho_n)\,h}\,\mathbf y_n,
\label{eq:expprop}
\end{equation}
which is unconditionally stable and exact for piecewise-constant $\rho$ (the
``analytic-within-a-step'' / PCA method). For time-varying $\rho$ the
piecewise-constant approximation incurs a local error $\mathcal O(h^2\dot\rho)$.
The action $e^{\mathbf M h}\mathbf v$ need not be formed densely; Krylov-subspace
methods compute it in $\mathcal O(m)$ matvecs via
$e^{\mathbf M h}\mathbf v\approx\|\mathbf v\|\,V_m e^{H_m h}e_1$, where $(V_m,H_m)$
is the Arnoldi factorisation of $\mathbf M$~\cite{molervanloan}.

For the \emph{coupled} neutronic--thermal system, Strang operator splitting advances
the neutronic flow $\Phi^N$ and thermal flow $\Phi^T$ as
\begin{equation}
\mathbf y_{n+1}=\Phi^N_{h/2}\circ\Phi^T_{h}\circ\Phi^N_{h/2}\,\mathbf y_n,
\end{equation}
with global error $\mathcal O(h^2)$ provided each sub-flow is solved at least to
second order; the splitting error is governed by the commutator $[\mathbf
M_N,\mathbf M_T]$ and vanishes when the blocks commute.

\section{Error bounds via the logarithmic norm}
\label{sec:lognorm}

Classical Lipschitz error bounds are useless for stiff systems (the Lipschitz
constant is $\sim|s_G|$). The sharp tool is the \emph{logarithmic norm}
$\mu(\mathbf M)=\lim_{h\to0^+}(\|\mathbf I+h\mathbf M\|-1)/h$; in the $2$-norm
$\mu_2(\mathbf M)=\lambda_{\max}\!\big(\tfrac12(\mathbf M+\mathbf M^{\!\top})\big)$.

\begin{theorem}[One-sided error bound]
\label{thm:lognorm}
If the exact and numerical flows satisfy $\dot{\mathbf x}=\mathbf M\mathbf x$ with
local truncation error $\boldsymbol\tau_n$, the global error
$\mathbf e_n=\mathbf x_n-\mathbf x(t_n)$ obeys
\begin{equation}
\|\mathbf e_n\|\le e^{\mu(\mathbf M)\,t_n}\|\mathbf e_0\|
+\sum_{k=1}^{n}e^{\mu(\mathbf M)(t_n-t_k)}\|\boldsymbol\tau_k\|.
\end{equation}
\end{theorem}
\begin{proof}[Sketch]
From $\dot{\mathbf e}=\mathbf M\mathbf e+\boldsymbol\tau$, the differential
inequality $\dfrac{\dd}{\dd t}\|\mathbf e\|\le\mu(\mathbf M)\|\mathbf
e\|+\|\boldsymbol\tau\|$ holds because
$\dfrac{\dd}{\dd t}\|\mathbf e\|\le\mu(\mathbf M)\|\mathbf e\|$ for the homogeneous
part; Gr\"onwall's inequality then yields the bound.
\end{proof}

The bound is informative precisely when $\mu(\mathbf M)$ is small or negative
(stable, damped dynamics) even though $\|\mathbf M\|$ is enormous; for a
prompt-critical excursion $\mu(\mathbf M)>0$ and the bound correctly predicts
exponential error growth---a quantitative warning that fast excursions demand
tight tolerances. This B-convergence framework~\cite{hairerwanner} is the modern
foundation for trusting stiff reactor-kinetics solvers.


\section*{Exercises}
\addcontentsline{toc}{section}{Exercises}
\begin{enumerate}[leftmargin=*,itemsep=2pt]
\item Why must a point-kinetics solver use a stiff integration method?
\item Describe how to validate a dynamics code against the inhour equation.
\item Explain what changes in the model output when the feedback coefficient sign is flipped.
\item What are the limitations of a lumped point-kinetics model for analysing a real, large reactor?
\item Outline what a higher-fidelity coupled neutronics/thermal-hydraulics code adds over a point model.
\end{enumerate}
